\section{Compute and Resource Markets}
\label{sec:markets}

The Lux provider model organizes compute, storage, and bandwidth as
continuous double-auction markets with on-chain settlement. Providers
post asks (offers to sell service at a given price). Users post bids
(willingness to pay for service at a given quality level). A matching
engine on the C-Chain clears the market each epoch.

\subsection{Market Structure}

Each resource type has an independent order book:

\begin{table}[H]
\centering
\caption{Resource market structure for the Lux decentralized cloud.}
\label{tab:markets}
\begin{tabular}{llll}
\toprule
\textbf{Market} & \textbf{Unit} & \textbf{Quality Dimensions} & \textbf{Settlement} \\
\midrule
Inference & GPU-second & Latency, throughput, model & Per-request \\
Storage & GB-month & Durability, availability, region & Per-epoch \\
Bandwidth & GB-transfer & Latency, throughput, region & Per-request \\
Recovery & Share-epoch & Availability, threshold & Per-epoch \\
Gateway & Request & Latency, concurrent connections & Per-request \\
\bottomrule
\end{tabular}
\end{table}

\subsection{Inference Market}

The inference market is the most complex, as it must match heterogeneous
compute capabilities (GPU type, VRAM, supported models) with heterogeneous
demand (model size, batch size, latency requirements).

An inference ask has the form:

\begin{equation}
    \mathsf{Ask}_{\text{infer}} = (\mathsf{gpu\_type}, \mathsf{vram}, \mathsf{models}[], \tau_{\max}, p_{\text{per\_token}})
\end{equation}

An inference bid has the form:

\begin{equation}
    \mathsf{Bid}_{\text{infer}} = (\mathsf{model}, \mathsf{max\_tokens}, \tau_{\max}, p_{\max}, \mathsf{mode})
\end{equation}

where $\mathsf{mode} \in \{\mathsf{TEE}, \mathsf{FHE}, \mathsf{plaintext}\}$
specifies the privacy level. TEE and FHE modes command a premium (typically
$1.3\times$ and $3\times$ respectively) due to the computational overhead.

\subsection{Market Clearing}

The matching engine runs a modified second-price auction each epoch.
For each resource type, bids and asks are sorted by price. Matches
are made greedily by descending bid price against ascending ask price.
The clearing price for each match is the ask price (second-price
semantics), incentivizing truthful bidding:

\begin{equation}
    p_{\text{clear}}(b, a) = p_a \quad \text{where } p_b \geq p_a
\end{equation}

\begin{theorem}[Truthful Bidding]
Under second-price semantics with on-chain settlement, truthful bidding
(reporting true valuation) is a weakly dominant strategy for users.
\end{theorem}

\begin{proof}
Standard Vickrey auction argument. If a user bids above true valuation
and wins, they may pay more than the resource is worth. If they bid below
and lose, they miss a profitable trade. Bidding true valuation maximizes
expected surplus in all cases.
\end{proof}

\subsection{Dynamic Pricing}

The base price for each resource type adjusts each epoch based on
utilization, following an EIP-1559-inspired mechanism:

\begin{equation}\label{eq:pricing}
    p_{t+1}^{\text{base}} = p_t^{\text{base}} \cdot \left(1 + \frac{1}{8} \cdot \frac{U_t - U_{\text{target}}}{U_{\text{target}}}\right)
\end{equation}

where $U_t$ is the utilization ratio (fraction of available capacity
consumed) and $U_{\text{target}} = 0.5$ is the target utilization.
When utilization exceeds 50\%, prices rise; when it falls below, prices
decrease. The $1/8$ dampening factor prevents price oscillation.

